\tocsection{\textbf{10/15/26 Lecture 14:} Random Processes and Moments}

If $(S, \cal{S})$ is a measurable space, then any $A \subseteq S$ (we don't need $A \in \cal{S}$!) becomes a measurable space with the $\sigma$-field $A \cap \cal{S} = \{A \cap B \col B \in \cal{S}\}$. Then, any random element of $(A, A \cap \cal{S})$ can be thought of as a random element in $S$.\\

Conversely, if $\xi$ is a random element in $S$ such that $\xi \in A$ a.s. (that is, $P(\xi^{-1}(A^{c})) = 0$) for some $A \in \cal{S}$, then $\xi = \eta$ almost surely for some random element $\eta$ in $A$.

\begin{definition}[\textbf{Definition 14.1:} Evaluation Maps]
    Fixing a measurable space $(S, \cal{S})$ and an abstract index set $T$, write $S^{T}$ for the class of functions $f : T \to S$, and let $\cal{S}^{T}$ denote the $\sigma$-field in $S^{T}$ generated by all \emph{evaluation maps} $\pi_t : S^{T} \to S$, $t \in T$, given by
    \[ \pi_t f = f(t) \]
\end{definition}

If $X : \Omega \to \cal{U} \subseteq S^{T}$, then clearly $X_t = \pi_t \circ X$ maps $\Omega$ into $S$. Thus, $X$ may be regarded as a function $X(\omega, t) = X_t(\omega)$ from $\Omega \times T$ into $S$. That is,
\[ X(\omega, t) = (X(\omega))(t) \]

\begin{example}{\textbf{Example}}
    Let $S = \R$ and $T = [0, 1]$, so $S^{T} = \R^{[0, 1]}$ is the set of all functions $f : [0, 1] \to \R$. Let $\cal{U} = C([0, 1], \R)$ be the continuous functions from $[0, 1]$ to $\R$. Then, a map $X : \Omega \to C([0, 1], \R)$ can be thought of as a function $X : \Omega \times [0, 1] \to \R$ by
    \[ X(\omega, t) = (X(\omega))(t) \]
\end{example}

\begin{lemma}[\textbf{Lemma 14.2}]
    Fix a measurable space $(S, \cal{S})$, an index set $T$, and a subset $\cal{U} \subseteq S^{T}$. Then, a function $X : \Omega \to \cal{U}$ is $\cal{U} \cap \cal{S}^{T}$-measurable if, and only if, $X_t : \Omega \to S$ is $\cal{S}$-measurable for all $t \in T$.
\end{lemma}

\begin{definition}[\textbf{Definition 14.3:} Random Processes]
    A mapping $X$ as above is called an \emph{$S$-valued random process on $T$ with paths in $\cal{U}$}. By the lemma, we can regard $X$ as a collection of random elements $X_t$ in the \emph{state space} $S$.
\end{definition}

\fig{0.4}{Lecture14Fig1}{}

\begin{definition}[\textbf{Definition 14.4:} Finite-Dimensional Distributions]
    For any random elements $\xi$ and $\eta$ in a common measurable space, the equality $\xi \overset{d}{=} \eta$ means that $\xi$ and $\eta$ have the same distribution, or $\cal{L}(\xi) = \cal{L}(\eta)$. If $X$ is a random process on some index set $T$, the associated \emph{finite-dimensional distributions} are given by
    \[ \mu_{t_1, \ldots, t_n} = \cal{L}(X_{t_1}, \ldots, X_{t_n}), \quad t_1, \ldots, t_n \in T,\ n \in \N \]
\end{definition}

\fig{0.45}{Lecture14Fig2}{}

\begin{lemma}[\textbf{Proposition 14.5:} Finite-Dimensional Distributions]
    Fix $S$, $T$, and $\cal{U} \subseteq S^{T}$, and let $X$ and $Y$ be processes on $T$ with paths in $\cal{U}$. Then, $X \overset{d}{=} Y$ if, and only if,
    % FIXED: page has X_{t_0} as the first entry; changed to X_{t_1} to match t_1, ..., t_n.
    \[ (X_{t_1}, \ldots, X_{t_n}) \overset{d}{=} (Y_{t_1}, \ldots, Y_{t_n}), \quad t_1, \ldots, t_n \in T,\ n \in \N \]
\end{lemma}
\begin{proof}{Proof:}
    % green text omitted here -- add your own
\end{proof}

\begin{definition}[\textbf{Definition 14.6:} Distribution Functions]
    For any random vector $\xi = (\xi_1, \ldots, \xi_d)$ in $\R^{d}$, we define the associated \emph{distribution function} $F$ by
    \[ F(x_1, \ldots, x_d) = P\left(\bigcap_{k \leq d} \{\xi_k \leq x_k\}\right), \quad x_1, \ldots, x_d \in \R \]
\end{definition}

\begin{lemma}[\textbf{Lemma 14.7:} Distribution Functions]
    Let $\xi$ and $\eta$ be random vectors in $\R^{d}$ with distribution functions $F$ and $G$. Then, $\xi \overset{d}{=} \eta$ if, and only if, $F = G$.
\end{lemma}
\begin{proof}{Proof:}
    This follows from the $\pi$-$\lambda$ theorem.
\end{proof}

\begin{definition}[\textbf{Definition 14.8:} Expected Value]
    The \emph{expected value}, \emph{expectation}, or \emph{mean} of a random variable $\xi$ is defined as
    \[ \mathbb{E}[\xi] = \int_{\Omega} \xi dP = \int_{\R} x (P \circ \xi^{-1})(dx) \]
    That is, we can either integrate $\xi$ over $\Omega$, or integrate the possible images with respect to the push-forward measure.
\end{definition}

For any random element $\xi$ in some measurable space $S$ and an arbitrary measurable function $f : S \to \R$,
\begin{align*}
    \mathbb{E}[f(\xi)] &= \int_{\Omega} f(\xi) dP = \int_{S} f(s)(P \circ \xi^{-1})(ds)\\
    &= \int_{\R} x (P \circ (f \circ \xi)^{-1})(dx)
\end{align*}
provided one of these three integrals exists. Integrals over a measurable subset $A \subseteq \Omega$ are often denoted by
\[ \mathbb{E}[\xi; A] = \mathbb{E}[\xi 1_{A}] = \int_{A} \xi dP, \quad A \in \cal{A} \]

\begin{definition}[\textbf{Definition 14.9:} Absolute Moments]
    For any random variable $\xi$ and constant $p > 0$, the integral $\mathbb{E}\abs{\xi}^{p} = \norm{\xi}_{p}^{p}$ is called the \emph{$p$th absolute moment} of $\xi$.
\end{definition}

By H\"older's inequality (or Jensen's, using the convexity of $x^{p}$), we see that $\norm{\xi}_{p} \leq \norm{\xi}_{q}$ for $p \leq q$, so the corresponding $L^{p}$-spaces are non-increasing in $p$.\\

Here is why. We want to show that $(\mathbb{E}\abs{\xi}^{p})^{1/p} \leq (\mathbb{E}\abs{\xi}^{q})^{1/q}$. Let $\eta = \abs{\xi}^{p}$. So, we want to show that $\mathbb{E}[\abs{\eta}] \leq \mathbb{E}[\abs{\eta}^{q/p}]^{p/q}$, or $\mathbb{E}[\abs{\eta}] \leq \mathbb{E}[\abs{\eta}^{r}]^{1/r}$ where $r = q/p \geq 1$. Then, by H\"older's,
\[ \mathbb{E}[\abs{\eta}] = \mathbb{E}[\abs{\eta} \cdot 1] \leq \norm{\eta}_{r}\norm{1}_{s} = \norm{\eta}_{r} = \mathbb{E}[\abs{\eta}^{r}]^{1/r} \]
where $\norm{1}_{s} = 1$ as we have a probability space, and where we use $p \leq q$ to ensure $r = \frac{q}{p} \geq 1$, so that there exists $s$ with $\frac{1}{r} + \frac{1}{s} = 1$.\\

\begin{definition}[\textbf{Definition 14.10:} Moments]
    If $\xi \in L^{p}$ and either $p \in \N$ or $\xi \geq 0$, define the \emph{$p$th moment} of $\xi$ as $\mathbb{E}[\xi^{p}]$.
\end{definition}
Note that we require $\xi \geq 0$ or $p \in \N$ to ensure we don't raise negative numbers to rational powers. 


\begin{lemma}[\textbf{Lemma 14.11:} Moments and Tails]
    For any random variable $\xi \geq 0$,
    \[ \mathbb{E}[\xi^{p}] = p\int_{0}^{\infty} P\{\xi > t\}t^{p - 1}dt = p\int_{0}^{\infty} P\{\xi \geq t\}t^{p - 1}dt, \quad p > 0 \]
\end{lemma}
\begin{proof}{Proof:}
    We see that
    \begin{align*}
        \mathbb{E}[\xi^{p}] &= \int_{\Omega} \xi^{p}P(d\omega)\\
        &= \int_{\Omega}\left(\int_{0}^{\xi(\omega)} pt^{p - 1}dt\right)P(d\omega)\\
        &= \int_{\Omega}\left(\int_{0}^{\infty} pt^{p - 1}1\{t \leq \xi(\omega)\}dt\right)P(d\omega)\\
        &= \int_{0}^{\infty}\int_{\Omega} pt^{p - 1}1\{t \leq \xi(\omega)\}P(d\omega)dt\\
        &= \int_{0}^{\infty} pt^{p - 1}P(\xi \geq t)dt
    \end{align*}
    where we swap the integrals using Tonelli and the fact th integrant is nonnegative. The $>$ follows as the difference of these sets is measure zero. The $\zeta > t$ version follows as $P(\xi \geq t) = P(\xi > t)$ except at the countably many $t$ with $P(\xi = t) > 0$
    
    $\qed$
\end{proof}
